\section{Introduction}
\label{sec:1}
A frozen structure predictor already contains operators that transform residue
information into pair representations. An adapter can reuse this computation,
or learn a generic pair update from new sequence features. These alternatives
raise two distinct questions: whether native factor composition is a useful
constraint, and whether its output directions are compatible with the fixed
network that consumes them. Neither question is answered by parameter count alone.

We construct zero-initialized residuals around native operators and investigate
these questions with trained controls. Protenix-Mini provides the initial setting:
a residue-wise predictor uses frozen ESM2 features to modify native query factors,
while retaining the frozen decoder and downstream pair processing. A generic
pair head tests the overall parameterization. Fixed orthogonal output rotations
then retain the factor construction and decoder spectrum while changing its
orientation relative to the downstream model. Rotated adapters receive training;
they are not native adapters damaged only at inference.

The evidence has three levels. First, independently locked Protenix comparisons
establish an advantage over the original generic head and a separate direction
effect. Second, fixed adapters remain useful on complete proteins longer than
those used for adaptation. Third, experiments using AlphaFold2 weights in
OpenFold test a substantially different architecture and expose the conditional
nature of the effect. Native orientation is advantageous in some AF2 settings,
but native factors do not consistently outperform a generic head given explicit
factor access. We report these boundaries alongside the positive results.

Our contribution is an operator-centered adaptation construction and controlled
empirical evidence about its orientation, scope and competitive limits. We do
not introduce OPM itself, claim recovery of evolutionary covariance, or identify
a complete mechanism for folding generalization.

\subsection{Related work}\label{related-work}

PLM-based folding includes HelixFold-Single, ESMFold and the ESM-conditioned
Protenix-Mini route \citep{protenix,helixfold,esmfold}. Our question is the parameterization of an internal
update in a frozen predictor. Adapters and Side-Tuning establish learning small
modules around pretrained networks \citep{adapters,sidetuning}, while ReZero studies residual
initialization \citep{rezero}. PiSSA, VeRA and ReFT respectively examine pretrained weight
directions, fixed random adaptation matrices and representation interventions
\citep{pissa,vera,reft}. Native directions or activation adaptation are therefore not new in
general. We study input-dependent native factor anchors and trained,
spectrum-preserving direction interventions in structure prediction. OpenFold
provides a trainable AF2 implementation \citep{openfold}; ColabFold is an additional AF2
execution pipeline \citep{colabfold}, not a second independent architecture in our study.
